% Löb, M. H. (1976). Embedding First Order Predicate Logic in Fragments of
% Intuitionistic Logic. The Journal of Symbolic Logic, 41(4), 705–718.
% Transcribed from scan. Some mathematical expressions marked [?] where the
% scan was ambiguous.

\documentclass[12pt,a4paper]{article}

\usepackage{amsmath}
\usepackage{amssymb}
\usepackage{amsthm}
\usepackage[T1]{fontenc}
\usepackage{microtype}
\usepackage{geometry}
\usepackage{enumitem}

\geometry{margin=1.25in}

%% ---- Shortcuts -------------------------------------------------------
\newcommand{\fA}{\mathfrak{A}}
\newcommand{\fB}{\mathfrak{B}}
\newcommand{\fC}{\mathfrak{C}}
\newcommand{\fM}{\mathfrak{M}}
%% Universal quantifier in the style of the paper: (x)A(x)
%% Written inline as, e.g., $(\eta)\fA(\eta)$

%% ---- Theorem environments --------------------------------------------
\newtheorem{lemma}{Lemma}
\newtheorem{corollary}[lemma]{Corollary}
\newtheorem{theorem}[lemma]{Theorem}

%% ---- Inference rule (fraction style) ---------------------------------
%% Usage: \infrule{premise(s)}{conclusion}
%% For named rules use, e.g., $\dfrac{P}{C}\;{\to}R$
%% Two-premise rules: \dfrac{P_1 \quad P_2}{C}

\begin{document}

%% -----------------------------------------------------------------------
\begin{center}
  {\large\bfseries
    EMBEDDING FIRST ORDER PREDICATE LOGIC IN\\[4pt]
    FRAGMENTS OF INTUITIONISTIC LOGIC}\\[10pt]
  {\large M. H. L\"{O}B}\\[6pt]
  \textit{The Journal of Symbolic Logic},
  Vol.\,41, No.\,4 (Dec., 1976), pp.\,705--718
\end{center}

\bigskip

Some syntactically simple fragments of intuitionistic logic possess
considerably expressive power compared with their classical counterparts.

In particular, we consider in this paper intuitionistic second order
propositional logic (ISPL) a formalisation of which may be obtained by adding
to the intuitionistic propositional calculus quantifiers binding propositional
variables together with the usual quantifier rules and the axiom scheme
$(Ex)(\mathfrak{A} = x)$, where $\mathfrak{A}$ is a formula not containing
$x$.

The main purpose of this paper is to show that the classical first order
predicate calculus with identity can be (isomorphically) embedded in ISPL.

It turns out an immediate consequence of this that the classical first order
predicate calculus with identity can also be embedded in the fragment
$(\mathrm{PL}_A)$ of the intuitionistic first order predicate calculus whose
only logical symbols are $\to$ and $(.)$ (universal quantifier) and the only
nonlogical symbol (apart from individual variables and parentheses) a single
monadic predicate letter.

Another consequence is that the classical first order predicate calculus can
be embedded in the theory of Heyting algebras.

The undecidability of the formal systems under consideration evidently follows
immediately from the present results.

We shall indicate how the methods employed may be extended to show also that
the intuitionistic first order predicate calculus with identity can be embedded
in both ISPL and $\mathrm{PL}_A$.

For the purpose of the present paper it will be convenient to use the
following formalisation $(S)$ of ISPL based on \cite{Schuttif1950}, rather
than the one given above.

%% -----------------------------------------------------------------------
\bigskip
\noindent\textbf{The formal system $S$.}\quad
The symbols of $S$ are
\begin{enumerate}
\item free variables: $a, b, c, \ldots$,
\item bound variables: $x, y, z, \ldots$,
\item logical symbols: the symbol $\to$ for implication, and the symbol $(.)$
  for universal quantification, where the dot stands for a bound variable,
\item parentheses.
\end{enumerate}

\noindent\textbf{Formulae of $S$.}
\begin{enumerate}
\item any free variable is a formula,
\item if $\mathfrak{A}$ and $\mathfrak{B}$ are formulae, then so is
  $(\mathfrak{A} \to \mathfrak{B})$,
\item if $\mathfrak{A}(\xi)$ is a formula in which the free variable $\xi$
  occurs at certain places then $(\eta)\mathfrak{A}(\eta)$ is a formula,
  provided $\eta$ is a bound variable not occurring in $\mathfrak{A}(\xi)$.
\end{enumerate}

An expression obtained from a formula by replacing some or all occurrences of
free variables by bound variables will be called a \emph{quasi formula}.

To simplify the notation we shall omit outermost parentheses from formulae.
Furthermore we shall write
$\mathfrak{A}_1, \mathfrak{A}_2, \cdots, \mathfrak{A}_n \to \mathfrak{B}$
for the formula
$\mathfrak{A}_1 \to (\mathfrak{A}_2 \to \cdots \to (\mathfrak{A}_n \to
\mathfrak{B})\cdots)$.

The formulae $\mathfrak{A}_1, \mathfrak{A}_2, \cdots, \mathfrak{A}_{n-1}$
will be said to be the \emph{antecedents} and $\mathfrak{B}$ the
\emph{succedent} of the formula
$\mathfrak{A}_1, \mathfrak{A}_2, \cdots, \mathfrak{A}_n \to \mathfrak{B}$.

Capital Greek letters will be used to denote sequences of formulae occurring
as antecedents.

In a formula of the form $\Gamma \to \mathfrak{B}$, $\Gamma$ may also be
empty. If so the expression $\Gamma \to \mathfrak{B}$ stands for
$\mathfrak{B}$. If $\mathfrak{A}$ is a formula of the form
$\Gamma \to \mathfrak{B}$, and $\mathfrak{B}$ is a formula having no
antecedents, then $\mathfrak{B}$ will be said to be the \emph{end formula}
of $\mathfrak{A}$.

\medskip
\noindent\textbf{Axioms of $S$.}\quad
All formulae of the form $\xi \to \xi$, where $\xi$ is a free variable.

\medskip
\noindent\textbf{Inferences of $S$.}\quad
They fall into two classes.

\smallskip
\noindent\textit{I.\ Structural inferences.}
\[
\text{Permutation}\quad
\dfrac{\Gamma,\,\mathfrak{A},\,\mathfrak{B},\,\Delta \to \mathfrak{C}}
     {\Gamma,\,\mathfrak{B},\,\mathfrak{A},\,\Delta \to \mathfrak{C}}
\qquad\qquad
\text{Contraction}\quad
\dfrac{\mathfrak{A},\,\mathfrak{A} \to \mathfrak{C}}
      {\mathfrak{A} \to \mathfrak{C}}
\]

\noindent\textit{II.\ Genetic inferences.}
\[
{\to}R\quad
\dfrac{\mathfrak{C}}{\mathfrak{A} \to \mathfrak{C}}
\qquad\qquad
{\to}L\quad
\dfrac{\Gamma \to \mathfrak{A}\qquad\mathfrak{B} \to \mathfrak{C}}
      {\Gamma,\,(\mathfrak{A} \to \mathfrak{B}) \to \mathfrak{C}}
\]
\[
{(.)R}\quad
\dfrac{\Gamma \to \mathfrak{A}(\xi)}{\Gamma \to (\eta)\mathfrak{A}(\eta)}
\qquad\qquad
{(.)L}\quad
\dfrac{\mathfrak{A}(\mathfrak{B}) \to \mathfrak{C}}
      {(\eta)\mathfrak{A}(\eta) \to \mathfrak{C}}
\]

The formula(e) above the stroke in an inference is (are) called
\emph{premise(s)}, and the formula below the stroke is called
\emph{conclusion} of the respective inference.

In $(.)R$ the free variable $\xi$ must not occur in the conclusion.

In the listed inferences formulae occurring as antecedents other than as
members of $\Gamma$ or as succedents other than those designated by
$\mathfrak{C}$ will be called \emph{principal formulae}.

Derivations in $S$ are trees of formulae starting from axioms and leading
downward by inferences to the lowest formula. The latter will be said to be
\emph{derivable in $S$}.

By the \emph{length} of a branch of a derivation we shall mean the number of
formulae on the branch. The length of the longest branch will be called the
\emph{length of the derivation}.

%% -----------------------------------------------------------------------
\medskip
A formula $\mathfrak{A}$ will be said to be \emph{derivably equivalent} to a
formula $\mathfrak{B}$ in $S$ if both $\mathfrak{A} \to \mathfrak{B}$ and
$\mathfrak{B} \to \mathfrak{A}$ are derivable in $S$.

Any two quasi formulae $\mathfrak{A}_x$ and $\mathfrak{B}_x$, obtained from
the formulae $\mathfrak{A}$ and $\mathfrak{B}$, respectively, by the same
substitution of bound variables (not occurring in $\mathfrak{A}$ and
$\mathfrak{B}$) for free variables will be said to be \emph{derivably
equivalent} if $\mathfrak{A}$ and $\mathfrak{B}$ are derivably equivalent.

An important property of $S$ proved first in \cite{Prawitz1965} is that the
cut
\[
\dfrac{\Gamma \to \mathfrak{A}\qquad\mathfrak{A} \to \mathfrak{C}}
      {\Gamma \to \mathfrak{C}}
\]
is a derived rule, i.e.\ that if the premises $\Gamma \to \mathfrak{A}$ and
$\mathfrak{A} \to \mathfrak{C}$ are derivable in $S$, then so is the
conclusion $\Gamma \to \mathfrak{C}$.

From this and \cite[p.\,67]{Prawitz1965} it follows in particular that $S$ is
equivalent to the formalisation of ISPL given above. It also follows easily
that the Hilbert type formalisation of ISPL is equivalent with $S$: the axioms
are all formulae of the form
\begin{align*}
&\mathfrak{A} \to (\mathfrak{B} \to \mathfrak{A}),\\
&(\mathfrak{A} \to (\mathfrak{B} \to \mathfrak{C}))
  \to (\mathfrak{A} \to \mathfrak{B}) \to (\mathfrak{A} \to \mathfrak{C}),\\
&(\xi)\mathfrak{A}(\xi) \to \mathfrak{A}(\mathfrak{B}).
\end{align*}
The inference rules are $(.)R$ and modus ponens
\[
\dfrac{\mathfrak{A}\qquad\mathfrak{A} \to \mathfrak{B}}{\mathfrak{B}}.
\]

%% -----------------------------------------------------------------------
\begin{lemma}
Let $\mathfrak{A}^{[\xi]}$ be the formula obtained from $\mathfrak{A}$ by
replacing all occurrences of the free variable $\xi$ by the formula
$\xi \to \xi$. For any formula $\mathfrak{A}$, the pair of formulae
$\xi \to \mathfrak{A}$ and $\xi \to \mathfrak{A}^{[\xi]}$ are derivably
equivalent in $S$.
\end{lemma}

\begin{proof}
By induction on the complexity of $\mathfrak{A}$. (The complexity of formulae
being the number of logical symbols occurring in it.)
\end{proof}

\begin{corollary}
Let $\mathfrak{B}$ be a formula of the form
$\mathfrak{A}_1, \mathfrak{A}_2, \cdots, \mathfrak{A}_n \to a$, where
$\mathfrak{A}_i$ is of the form $\xi \to \mathfrak{B}_i$ ($\xi$ being a free
variable). Then if $\mathfrak{B}^{[\xi]}$ is derivable in $S$ so is
$\mathfrak{B}$.
\end{corollary}

\begin{proof}
Suppose $\mathfrak{B}^{[\xi]}$ is derivable in $S$, then so is the formula
$\mathfrak{B}_1^{[\xi]}, \cdots, \mathfrak{B}_n^{[\xi]}$ and hence also the
formula
$(\xi \to \mathfrak{B}_1^{[\xi]}), \cdots, (\xi \to \mathfrak{B}_n^{[\xi]})
 \to (\xi \to \mathfrak{B}_{n+1}^{[\xi]})$.
Then by Lemma~1, the formula
$(\xi \to \mathfrak{B}_1), \cdots, (\xi \to \mathfrak{B}_n)
 \to (\xi \to \mathfrak{B}_{n+1})$, i.e.\ $\mathfrak{B}$ is derivable in $S$.
\end{proof}

%% -----------------------------------------------------------------------
By \emph{parameters} we shall mean certain distinguished free variables (to be
specified in the context whenever necessary).

A set $M$ of formulae will be said to be \emph{closed with respect to
substitution of formulae} (c.s.f.), if whenever $\mathfrak{A}(\xi)$ is a
formula of $M$ with $\xi$ a free variable other than a parameter,
$\mathfrak{A}(\mathfrak{B})$ is also a formula of $M$ for any formula
$\mathfrak{B}$ (provided $\mathfrak{A}(\mathfrak{B})$ is a formula).

Let $M$ be a set of formulae of $S$ and $\xi$ a free variable. We define
$C(M, \xi)$ to be the smallest closed set of primary formulae satisfying the
following conditions.
\begin{enumerate}[(i)]
\item $\xi \in C(M, \xi)$,
\item if $\mathfrak{A}$ and $\mathfrak{B}$ belong to $C(M, \xi)$, then
  $\mathfrak{A} \to \mathfrak{B}$ is in $C(M, \xi)$,
\item for each formula $\mathfrak{B}$, $\mathfrak{A}(\mathfrak{B})$ is in
  $C(M, \xi)$ (for bound variables $\zeta$ such that
  $(\zeta)\mathfrak{A}(\zeta)$ is a formula).
\end{enumerate}

A formula $\mathfrak{A}$ will be said to be an \emph{$M$-formula} if
$\mathfrak{A}$ and all formulae in $\Gamma$ belong to $M$.

To be a \emph{$(M,\xi)$-formula}, if $\mathfrak{A}$ is either an
$M$-formula or belongs to $C(M, \xi)$.

A tree of $(M,\xi)$-formulae which is constructed like a derivation in $S$
except that any formula at the top of any branch is either an $M$-formula
derivable in $S$ or a formula belonging to $C(M, \xi)$ of the form
$\xi \to \mathfrak{A}$ will be called an \emph{$(M,\xi)$-derivation}.

%% -----------------------------------------------------------------------
\begin{lemma}
For any set $M$ of formulae and free variable $\xi$, if $\mathfrak{A}$ is a
formula belonging to $C(M, \xi)$, then $\mathfrak{A}^{[\xi]}$ is derivable
in $S$.
\end{lemma}

\begin{proof}
By induction with respect to the complexity of $\mathfrak{A}$.
If the length of the derivation of $\mathfrak{A}$ is $0$, then $\mathfrak{A}$
is an axiom. Since $\mathfrak{A}$ belongs to $C(M,\xi)$ it must have the form
$\xi \to \xi$. In $C(M,\xi)$ it must have the form $\xi \to \mathfrak{A}$.
\end{proof}

%% -----------------------------------------------------------------------
\begin{lemma}
If for a given parameter $\xi$, $M$ is a set of formulae c.s.f.\ of the form
$\xi \to \mathfrak{A}$, then any formula $\mathfrak{A}$ of $C(M, \xi)$ which
is derivable in $S$ has an $(M, \xi)$-derivation.
\end{lemma}

\begin{proof}
By induction with respect to the length of the derivation in $S$ of
$\mathfrak{A}$.

Suppose now that the length of the derivation of $\mathfrak{A}$ is $n+1$ and
that the lemma holds for all formulae of $C(M,\xi)$ with derivation length
$n$.

It is easily verified by inspection of the inferences of $S$ that all
inferences except $\to L$ have the property that if the conclusion belongs to
$C(M,\xi)$ then so does the premise(s). Hence the induction hypothesis holds
for the premises. Also for the conclusion.

We now consider the case when the last inference in the derivation of
$\mathfrak{A}$ is a $\to L$ of the form

\[
\dfrac{\Gamma \to \mathfrak{A}\qquad\mathfrak{B} \to \mathfrak{C}}
      {\Gamma,\,(\mathfrak{A} \to \mathfrak{B}) \to \mathfrak{C}}
\]

(a) If both $\mathfrak{A}$ and $\mathfrak{B}$ belong to $C(M,\xi)$ then both
premises belong to $C(M,\xi)$. By the induction hypothesis the lemma holds for
the premises and must therefore also hold for the conclusion.

(b) Suppose $\mathfrak{A}$ belongs to $M$ and $\mathfrak{B}$ to $C(M,\xi)$.
Let $\mathfrak{C}_1, \cdots, \mathfrak{C}_n$ be a list (possibly empty) of
all formulae belonging to $M$ in the derivation, and then belonging to
$C(M,\xi)$.
Now consider the formula $\Gamma^{[k]}$ which is the sequence obtained from
$\Gamma$ by replacing each formula $\mathfrak{B}$ by $\mathfrak{B}^{[k]}$.
Clearly this formula is derivable in $S$. Let $\Gamma_n$ be the sequence ...
[derivation argument continues]

Thus since $\mathfrak{A}$ is a formula belonging to $F_n$, $\tau(\mathfrak{A})^{[k]}$
is derivable in $S$.

Now let $S^-$ be the formal system obtained from $S$ by restricting
$\mathfrak{B}$ in $(.)L$ to free variables. Since $\sigma^*$ faithfully maps
any derivation of $\tilde{\mathfrak{A}}$ in $\mathrm{CPL}^-$ step by step
into an $(M, a)$-derivation, say, of $\sigma^*(\mathfrak{A})$ any $(.)I$.

\[
\dfrac{\mathfrak{A}(\mathfrak{B}) \to \mathfrak{C}}{\mathfrak{A}(\mathfrak{C})}
\qquad
\dfrac{(\xi)\mathfrak{A}(\xi) \to \mathfrak{C}}{\mathfrak{A}}
\]

in $D$ is such that $\mathfrak{B}$ is a free variable. Hence $D$ is a
derivation in $S^-$.
\end{proof}

%% -----------------------------------------------------------------------
\noindent
Thus the theorem remains true if $S$ is replaced by $S^-$. In fact, the
theorem still remains true if $S$ is replaced by $S^-$ and $\sigma$, where
$\tilde{\sigma}$ is defined like $\sigma^*$ except for taking
$\neg I(\xi)$ in place of $U(\xi)$. Clearly if $\tilde{\mathfrak{A}}$ is
derivable in $\mathrm{CPL}^-$ then $\tilde{\sigma}(\tilde{\mathfrak{A}})$ is
derivable in $S^-$. On the other hand, if $\tilde{\sigma}(\tilde{\mathfrak{A}})$
is derivable in $S^-$ it is

\bigskip
\noindent
We now add to the $(M,a)$-derivation of the right premise the inferences

\[
\dfrac{\Gamma,\,\mathfrak{B},\,\mathfrak{B}_{1s} \to \mathfrak{C}}
      {\Gamma,\,\mathfrak{B} \to \mathfrak{B}_{1s} \to \mathfrak{C}}
\qquad
\dfrac{\Gamma \to \mathfrak{C}\qquad\mathfrak{B} \to \mathfrak{D}}
      {\Gamma,\,(\mathfrak{B} \to \mathfrak{C}) \to \mathfrak{D}}
\]

%% -----------------------------------------------------------------------
\begin{corollary}
Let $M$ be as in Lemma~3 and $\mathfrak{A}$ a formula in $C(M,\xi)$
derivable in $S$ of the form $\Gamma \to \xi$, where $\Gamma$ is a sequence
of formulae $\mathfrak{A}_1, \cdots, \mathfrak{A}_n$, such that
$\mathfrak{A}_i$ either belongs to $M$ or is of the form
$\mathfrak{B}_i \to \xi$ for some $\mathfrak{B}_i$, in $M$. Then the
formulae of the form $\Delta \to \mathfrak{B}_j$, for some $j$
$(1 \le j \le n)$, where $\Delta$ is a subsequence of formulae of $\Gamma$
belonging to $M$, are derivable in $S$.
\end{corollary}

\begin{corollary}
Let $M$ be as in Lemma~3 and $\mathfrak{A}$ a formula in $C(M,\xi)$
derivable in $S$ of the form $\Gamma \to \xi$, where $\Gamma$ is a sequence
$\mathfrak{B}_1, \cdots, \mathfrak{B}_{l0} \to \xi$ $(1 \le j \le m)$
where $\mathfrak{B}_{1j}$ belongs to $M$. Then the formulae of the form
$\mathfrak{B}_{1,1}, \cdots, \mathfrak{B}_{l,n} \to \mathfrak{B}_{0}$ and
all $1 \le l \le j$, are derivable in $S$, where $\Delta$ is a subsequence
of formulae of $\Gamma$ belonging to $M$.
\end{corollary}

\begin{proof}
By the lemma $\mathfrak{A}$ has an $(M,\xi)$-derivation, $D$. The corollary
follows by a simple induction with respect to the length of $D$ (i.e.\ the
length of a maximal branch of $D$).
\end{proof}

%% -----------------------------------------------------------------------
\noindent
Let $M'$ be the set of formulae such that for any set of formulae $ξ \to \mathfrak{C}$,
belonging to $M$, the formula $\xi \to \mathfrak{C}, \&\, \xi \to \mathfrak{C}\, \& \,
\cdots$ belongs to $M'$.

Then putting $\mathfrak{B}_{1j} = \xi \to \mathfrak{D}_{1j}$,
$\mathfrak{B}_{1s}, \cdots, \mathfrak{B}_{1n} \to \xi$ are derivably
equivalent with $(\xi \to \mathfrak{D}_{1s}) \& \cdots \& (\xi \to \mathfrak{D}_{1n})
\to \xi$. Then from the previous corollary it follows easily.

Let $\Gamma^*$ be the sequence of formulae obtained from $\Gamma$ by replacing
each formula $\mathfrak{B}$ by $(\mathfrak{B} \to \mathfrak{D}_{1s}) \&
\cdots \& (\mathfrak{B} \to \mathfrak{D}_{1n}) \to \mathfrak{B}$. Then clearly
$\Gamma^* \to \xi$ is derivable in $S$. By Corollary (i) a formula of the
form $\Theta \to \xi$ where $\Theta$ is a subsequence of formulae of
$M$ belonging to $\Gamma^*$

%% -----------------------------------------------------------------------
\bigskip
\noindent
The language $L$ (the negative fragment of) intuitionistic predicate logic PL
is obtained from the language of $S$ by adding predicate letters and function
symbols of suitable finite types%
\footnote{It should be noted that the standard results for the first order
predicate calculus (e.g.\ completeness and cut elimination) continue to hold
if function symbols of higher types occur in terms, provided that all terms
have the type of individuals.}
and $\bot$ (absurdity) as a logical symbol. Terms, prime formulae and formulae
of $L$ may then be defined in the usual way. The rules of PL are the same as
those for $S$, except for adding $\bot$ in $(.)L$ to be a term instead of a
formula,%
\footnote{From this observation it follows easily that $\mathrm{CPL}^*$ can be
embedded in the theory of positive implicational algebras (PIA), i.e.\ the
classical theory based on the axioms $\alpha$(i), (ii), (v). To see this let
$\tilde{\mathfrak{A}}$ be the formula obtained from $\tilde{\mathfrak{A}}$,
(here taking $U(\xi)$ to be $\neg I(\xi)$) by replacing quasi prime formulae
of the form $I(\xi)$ by $(I(\xi) \to c_i) \to c_i$:
\[
I(c \to \xi) \supset c_i) \to ((c \to \xi_2) \supset c_i) \to \cdots
\to ((c \to \xi_n) \supset c_i) \supset b).
\]
Now if $\tilde{\mathfrak{A}}$ is derivable in $\mathrm{CPL}^-$, then so is
$\tilde{\mathfrak{A}}$, and hence also $\tilde{\mathfrak{A}}$. Since $\tilde{\mathfrak{A}}$
has a derivation, $D$ say, in PIA. Then by dropping all occurrences of $I$
in $D$ and replacing $\supset$ by $\to$ throughout we obtain a derivation of
$\sigma(\mathfrak{A}_{11})$ in $S^-$. Hence $\tilde{\mathfrak{A}}$ is derivable
in $S^-$.

Moreover, since the theory of Heyting algebras is a conservative extension of
PIA we have that $\mathrm{CPL}^*$ can be embedded in the theory of Heyting
algebras.}
and taking for axioms all formulae of the form $\mathfrak{A} \to \mathfrak{A}$
and $\bot \to \mathfrak{A}$.

A \emph{derivation in PL (CPL) from a class of formulae $K$}, if there is a
tree which is constructed like a derivation of PL (CPL), except for allowing
at the top of branches besides axioms to occur formulae of $K$ (which in the
case of CPL are classical).

A formula $\mathfrak{B}$ is \emph{derivable in PL (CPL) from a class of
formulae $K$}, if there is a tree which is constructed like a derivation of PL
(CPL) for allowing at the top of any branch to occur formulae of $K$.

A formula $\mathfrak{A}$ of $L$ will be said to be a \emph{primary formula},
if it is of the form $\mathfrak{A}_1, \cdots, \mathfrak{A}_n \to \mathfrak{A}_{n+1}$,
where each $\mathfrak{A}_i$ is a prime formula.

A set $N$ of primary formulae will be said to be \emph{closed} if it
satisfies the following conditions:
\begin{enumerate}[(i)]
\item if $\mathfrak{C}(\xi)$ belongs to $N$, then so does $\mathfrak{C}(\eta)$,
  for any free variable $\xi$ and term $\eta$ (where occurrences of $\xi$ in
  $\mathfrak{C}(\eta)$ occur in $\eta$ only),
\item if $\Gamma, \mathfrak{A}, \Delta \to \mathfrak{C}$ belongs to $N$ and
  $\mathfrak{A} \to \mathfrak{B}$ belongs to $N$, then so does
  $\Gamma, \mathfrak{B}, \Delta \to \mathfrak{C}$,
\item if $\mathfrak{A}, \mathfrak{A} \to \mathfrak{B}$ belongs to $N$, then
  so does $\mathfrak{A} \to \mathfrak{B}$,
\item if $\Gamma \to \mathfrak{B}$ and $\mathfrak{B} \to \mathfrak{C}$ belong
  to $N$, then so does $\Gamma \to \mathfrak{C}$.
\end{enumerate}
For any set $P$ of primary formulae the closure $\mathrm{C1}(P)$ will be the
smallest closed set of primary formulae containing $P$.

Let $\mathrm{PL}^*$ be the formal system obtained from PL by adding the cut
rule. [It] is easily seen that for a primary formula $\mathfrak{A}$,
$\mathfrak{A}$ is derivable from $\mathrm{C1}(P)$ in PL if and only if
$\mathfrak{A}$ is derivable in $\mathrm{PL}^*$ from the closure of the set of
primary formulae $\bar{K}$, where $\bar{K}$ is the class of formulae
$\bar{\mathfrak{A}}$ with $\mathfrak{A}$ in $K$.

%% -----------------------------------------------------------------------
\begin{lemma}
Let $P$ be a set of primary formulae, then any formula derivable from $P$ in
$\mathrm{PL}^*$ is derivable in PL from $\mathrm{C1}(P)$.
\end{lemma}

The proof is analogous to the proof in \cite[§3]{Schuttif1950} of the
corresponding result for classical predicate logic.

\begin{corollary}
(i) Let $\bar{\mathfrak{A}}$ be the formula obtained from $\mathfrak{A}$ by
replacing every quasi prime formula $(\mathfrak{B} \to \bot)$ in $\mathfrak{A}$
by $(\mathfrak{B} \to \bot) \to \bot$. Let $\bar{\Gamma}$ be a finite
sequence of the universal closures of quasi prime formulae and $\bar{\Gamma}$
the sequence obtained from $\Gamma$ by replacing $\Gamma$ by $\bar{\mathfrak{B}}$.
Then for every classical formula $\mathfrak{A}$, $\bar{\Gamma}, \bar{\mathfrak{A}} \to \mathfrak{A}$
is derivable in PL if and only if $\bar{\Gamma} \to \mathfrak{A}$ is derivable
in PL.
\end{corollary}

\begin{proof}
Suppose $\bar{\Gamma} \to \mathfrak{A}$ is derivable in PL, then clearly so
is $\bar{\Gamma} \to \bar{\mathfrak{A}}$. But since $\mathfrak{A}$ and
$\bar{\mathfrak{A}}$ are derivably equivalent in PL, $\bar{\Gamma} \to
\mathfrak{A}$ is derivable in PL. Hence $\bar{\Gamma} \to \mathfrak{A}$ is
derivable in PL.

Suppose now $\bar{\Gamma} \to \mathfrak{A}$ is derivable in PL. Since it is
easily seen that for a primary formula $\mathfrak{A}$, $\mathfrak{A} \to
\mathfrak{A}$ is derivable in PL, we have for every formula $\mathfrak{B}$ in
$\bar{\Gamma}$, $\mathfrak{B} \to \bar{\mathfrak{B}}$ is derivable in PL.
Hence $\bar{\Gamma} \to \mathfrak{A}$ is derivable in PL.
\end{proof}

\begin{corollary}
(ii) Let $K$ be a class of primary formulae. Then a classical formula
$\mathfrak{A}$ is derivable in PL from $K$ if and only if $\mathfrak{A}$
is derivable in CPL from $\bar{K}$, where $\bar{K}$ is the class of formulae
$\bar{\mathfrak{A}}$ with $\mathfrak{A}$ in $K$.
\end{corollary}

This follows immediately from Corollary (ii) of Lemma~4.

In particular, from Corollary (ii) of Lemma~4 it follows immediately that a
formula $\mathfrak{A}$ is derivable in classical predicate logic with identity
$\mathrm{CPL}^*$ if and only if $\mathfrak{A}$ is a classical formula of $L$
(possibly containing the identity symbol) derivable in PL from the closure of
the set of primary formulae
\begin{enumerate}[(i)]
\item $\xi = \xi$
\item $\xi = \eta \to (\mathfrak{A}(\xi) \to \mathfrak{A}(\eta))$, where
  $\mathfrak{A}(\xi)$ is a prime formula.
\end{enumerate}

A structure $\mathscr{S} = \langle D, I, \Rightarrow, \bigcap\rangle$ where
$D$ is a nonempty set, $I$ a subset of $D$, $\Rightarrow$ a mapping from
$D^2$ into $D$, and $\bigcap$ a mapping from $D^D$ into $D$, will be said to
be an \emph{$S$-structure} if it satisfies the following set $\alpha$ of
conditions:
\begin{enumerate}[(i)]
\item $I(\xi \Rightarrow (\eta \Rightarrow \xi))$,
\item $I((\xi \Rightarrow (\eta \Rightarrow \zeta)) \Rightarrow
    ((\xi \Rightarrow \eta) \Rightarrow (\xi \Rightarrow \zeta)))$,
\item $I(\bigcap \xi(\xi), \xi(\eta))$, where $\eta$ is a term,
\item if $I(\xi \Rightarrow \eta(\xi))$ belongs to $\alpha$, provided $\xi$
  is a free variable not occurring in the term
  $\xi \Rightarrow \bigcap_\eta \eta(x)$,
\item $I(\xi),\; I(\xi \Rightarrow \eta) \Rightarrow I(\eta)$.
\end{enumerate}
In (iii) and (v) $x$ is an arbitrary bound variable. In (i)--(v) the symbols
$I, \Rightarrow, \bigcap$ are used as formal symbols of $L$ with the
understanding that their interpretations in $\mathscr{S}$ are denoted by the
same symbol, respectively, and $\xi, \eta, \zeta$ are terms constructed from
free variables (taken to range over $D$) and the operation symbols $\Rightarrow$
and $\bigcap$.

The same convention will be adopted for the extensions of $S$-structures
considered below.

Evidently every $S$-structure is a lattice with respect to the ordering
relation $I(\xi \Rightarrow \eta)$. The lattice theoretic operations $\cap$
and $\cup$ may be defined such that for any two elements $\xi$ and $\eta$ of
$D$
\[
\xi \cap \eta \quad\text{stands for}\quad
\bigcap_{x} \bigl[(\xi \Rightarrow (\eta \Rightarrow x)) \Rightarrow x\bigr]
\]
and
\[
\xi \cup \eta \quad\text{stands for}\quad
\bigcap_{x} \bigl[(\xi \Rightarrow x) \cap (\eta \Rightarrow x)\bigr]
  \Rightarrow x.
\]
Moreover, any $S$-structure has a least element $0$ and a greatest element
$\mathbf{1}$ definable as $\bigcap_x x$ and $\bigcap_x (x \Rightarrow x)$,
respectively. Recalling the cut elimination property and Hilbert type
formalisation of ISPL it is easy to prove

%% -----------------------------------------------------------------------
\begin{lemma}
Let $\mathfrak{A}$ be a prime formula in the language of $S$-structures and
$\mathfrak{A}'$ the formula in the language of $S$ obtained from $\mathfrak{A}$
by dropping $I$ and replacing $\Rightarrow$ and $\bigcap_\xi$ by $\to$ and
$(\xi)$, respectively. Then $\mathfrak{A}$ belongs to $\mathrm{C1}(\alpha)$
if and only if $\mathfrak{A}'$ is derivable in $S$.
\end{lemma}

Note also that since all formulae of $\alpha$ are primary we can by Lemma~4
present the theory of $S$-structures in cut-free form taking $\mathrm{C1}(\alpha)$
as the set of nonlogical axioms.

We now extend the theory of $S$-structures by adding the parameters
$a, b, c_1, \cdots, c_n$ and an $n$-ary predicate letter $A$ subject to the
following set $\beta$ of conditions:
\begin{enumerate}[(vi)]
\item $I(a)$, $I(b)$, $I(c_1)$, $I(c_n)$,
\item $(I\xi \Rightarrow \eta)$,
  $A(\xi_1, \cdots, \xi_{n-1}, \xi, \xi_1, \cdots, \xi_n)
   \to A(\xi_1, \cdots, \xi_{n-1}, \eta, \xi_1, \cdots, \xi_n)$.
\end{enumerate}

Let $C_{LA}$ be the set of classical formulae derivable in PL from
$\mathrm{C1}(\alpha \cup \beta)$.

A structure $\langle D, I, A, \Rightarrow, \bigcap\rangle$ where
$A \subseteq D^n$ will be said to be an \emph{expanded $S$-structure} if it
is an $S$-structure if it satisfies the conditions $\beta$(vi)--(vii) under
an assignment of elements of $D$ to the parameters $a, b, c_1, \cdots, c_n$.

We now define a mapping $\tau$ of the set of formulae of the language of
$C_{LA}$ into the set of formulae of the language of $S$.
\begin{align*}
\tau(\xi) &= \xi, \quad\text{if $\xi$ is a variable,}\\
\tau(\xi = \eta) &= (\eta \Rightarrow \xi) \cap (\xi \Rightarrow \eta), \\
\tau(\bigcap_\xi(\xi)) &= (\eta)\,\tau(\xi(\eta)), \\
\tau(\bot) &= a, \\
\tau(\mathfrak{A} \mathbin{\&} \mathfrak{B}) &= \tau(\mathfrak{A}) \cap \tau(\mathfrak{B}), \\
\tau(\mathfrak{A} \vee \mathfrak{B}) &= \tau(\mathfrak{A}) \cup \tau(\mathfrak{B}), \\
\tau(\mathfrak{A} \to \mathfrak{B}) &= \tau(\mathfrak{A}) \Rightarrow \tau(\mathfrak{B}), \\
\tau((\xi)\mathfrak{A}(\xi)) &= (\xi)\,\tau(\mathfrak{A}(\xi)), \\
\tau((E\xi)\mathfrak{A}(\xi)) &= (E\xi)\,\tau(\mathfrak{A}(\xi)).
\end{align*}

In the proof of the next lemma reference will be made to the following sets of
formulae:
\begin{align*}
F_1 &\;=\;\text{the set of formulae of the form } b, c_1, \cdots, c_n \to \mathfrak{A},\\
F_2 &\;=\;\text{the set of formulae of the form } a \to \mathfrak{A}, \text{ where } \mathfrak{A} \text{ belongs to } F_3,\\
F_{3,i} &\;=\;\text{the set of formulae of the form } (c_1, \cdots, c_n \to \mathfrak{A}) \to c_i,\\
F_3  &\;=\;\text{the set of formulae of the form } h \to \mathfrak{A}, \text{ where } \mathfrak{A} \text{ belongs to } F_{3,i},\\
F_4  &\;=\;\text{the set of formulae of the form } \mathfrak{A}_i \cdots \mathfrak{A}_k \to \mathfrak{A}_{k+1}, \\
     &\phantom{\;=\;}\text{where } \mathfrak{A}_i \text{ belongs to } F_{4,n},\\
F_5  &\;=\;\text{the set of formulae of the form }
       a \to \mathfrak{A}, \text{ where } \mathfrak{A}_i \text{ belongs to } F_5,\\
F_6  &\;=\;\text{the set of formulae of the form } a \to \mathfrak{A}, \text{ where } \mathfrak{A} \text{ belongs to } F_5.
\end{align*}

%% -----------------------------------------------------------------------
Let $N$ be the union of $F_1$ and $F_4$, and $M$ the union of $F_2$ and $F_6$.

A quasi formula obtained from a formula $\mathfrak{B}$ in $M$ by substituting
bound variables not occurring in $\mathfrak{B}$ for one or more free variables
other than parameters will be said to be \emph{quasi formulae of $M$}.

Note that any (quasi) formula of $M$ has one of the following forms:
\begin{enumerate}[(i)]
\item $a, b, c_1, \cdots, c_n \to \mathfrak{A}$,
\item $a, b, [c_1, \cdots, c_n \to \mathfrak{A}_1], \cdots,
  [b, [c_1, \cdots, c_n \to \mathfrak{A}_m] \to c_n] \to b$.
\end{enumerate}
The (quasi) formulae $\mathfrak{A}_1$ in (i) and (ii) will be said to be
the main parts of the quasi formulae (i) and (ii), respectively.

An occurrence $\mathfrak{C}$ of a quasi formula will be said to be a main
part of an $(M, a)$-formula if $\mathfrak{C}$ is the main part of a quasi
subformula of $\mathfrak{A}$ belonging to $M$.

%% -----------------------------------------------------------------------
\begin{lemma}
For any formula $\mathfrak{A}$ of the language of $S$ there is at most one
formula $\mathfrak{B}$ of the language of $C_{LA}$ such that $\mathfrak{A} = \tau(\mathfrak{B})$.
\end{lemma}

This follows easily from the observation that any occurrence of $\to$ and
$(\xi)$ corresponds under $\tau$ to $\Rightarrow$ and $\bigcap$, respectively,
only if it is contained in a main part of $\mathfrak{A}$.

%% -----------------------------------------------------------------------
\begin{lemma}
A formula $\mathfrak{A}$ belongs to $C_{LA}$ if and only if $\tilde{\mathfrak{A}}$
is classical and $\tau(\tilde{\mathfrak{A}})$ is derivable in $S$.
\end{lemma}

\begin{proof}
Suppose $\mathfrak{A}$ belongs to $C_{LA}$. Then it has a derivation, $D$
say, in PL from $\mathrm{C1}(\alpha \cup \beta)$. Let $T(D)$ be the tree of
formulae obtained from $D$ by replacing each formula $\mathfrak{A}$ by
$\tau(\mathfrak{A})$. Then it is easily seen that each formula at the top of
$T(D)$ is derivable in $S$. Moreover, every inference of PL in $D$ is carried
into an inference of $S$. Hence $\tau(\mathfrak{A})$ is derivable in $S$.

Suppose now for classical $\mathfrak{A}$, $\tau(\mathfrak{A})$ is derivable
in $S$. Clearly $\tau^*(\mathfrak{A})$ is derivable in $S$. Suppose now
$\sigma^*(\mathfrak{A})$ is derivable in $S$. Then also the formula
$\mathfrak{B}$ obtained from $\mathfrak{A}$ by $\sigma^*$ by replacing
$c$ throughout by $a, b, c_1, \cdots, c_n$, and $b$ by $a$ is derivable in
$S$. But clearly $\mathfrak{B}$ is derivably equivalent with $\tau(\mathfrak{A})$.
Hence by the theorem, $\mathfrak{A}$ is derivable in $\mathrm{CPL}^*$.
\end{proof}

%% -----------------------------------------------------------------------
\noindent
Now let $S^-$ be the formal system obtained from $S$ by restricting
$\mathfrak{B}$ in $(.)L$ to free variables. Since $\sigma^*$ faithfully maps
any derivation of $\tilde{\mathfrak{A}}$ in $\mathrm{CPL}^*$ step by step
into an $(M, a)$-derivation, say, of $\sigma^*(\tilde{\mathfrak{A}})$ any
$(.)I$.

Now if $\tilde{\mathfrak{A}}$ is derivable in $\mathrm{CPL}^-$, then so is
$\tilde{\mathfrak{A}}$, and hence also $\tilde{\mathfrak{A}}$. Since
$\tilde{\mathfrak{A}}$ has a derivation, $D$ say, in PIA. Then by dropping
all occurrences of $I$ in $D$ and replacing $\supset$ by $\to$ throughout we
obtain a derivation of $\sigma(\mathfrak{A}_{11})$ in $S^-$. Hence
$\tilde{\mathfrak{A}}$ is derivable in $S^-$.

%% -----------------------------------------------------------------------
\bigskip
\noindent
We shall now prove more generally that if
$\mathfrak{A}_1, \cdots, \mathfrak{A}_n, A(\nu_1, \cdots, \nu_{k-1},
\nu_k, \nu_{k+1}, \cdots, \nu_n)$
is derivable in $S$ for arbitrary terms $\nu_1, \cdots, \nu_k, \nu_{k+1},
\nu_{k+1}, \cdots, \nu_n$ in the language of $S$-structures.
From these we get by cuts the formula
$I(\eta_h), \cdots, I(\eta_n), A(\nu_1, \cdots, \nu_{k-1}, \nu_k,
\nu_{k+1}, \cdots, \nu_n)$
for arbitrary terms $\nu_1, \cdots, \nu_k, \nu_{k+1}, \nu_n$ in the language
of $S$-structures.
From these we get by cuts the formula $A(\xi_1, \cdots, \xi_k, \xi_{k+1},
\cdots, \xi_n)$.
Hence $\mathfrak{A}$ belongs to $\mathrm{C1}(\alpha \cup \beta)$.

Hence the formulae $I(\eta_h), \cdots, I(\eta_n), A(\nu_1, \cdots,
\nu_{n-1}, \nu_k, \nu_{k+1}, \cdots, \nu_n)$ for any $\mathfrak{A}$ belonging
to $F_n$ $(t(\mathfrak{A}) = \tau(\tau_n))$ are derivable in $S$. Now
$\bar{\mathfrak{X}}_k = \tau(\xi_k)$, $\bar{\mathfrak{Y}}_h = \tau(\eta_h)$
and $\bar{\mathfrak{Z}}_k = \tau(\eta_k)$ for some terms $\xi_k$, $\eta_h$,
$\eta_k$ in the language of $S$-structures, and $\bar{\mathfrak{Z}}_k =
\tau(\xi_k)$.

%% -----------------------------------------------------------------------
\bigskip\noindent
Hence the formulae
\[
(c_k \to \mathfrak{B}_n^{[k]}), \cdots, (c_n \to \mathfrak{B}_n^{[k]}),
(c_k \to \bar{\mathfrak{X}}_k^{[k]}) \to (c_k \to \bar{\mathfrak{Z}}_k)
\quad (1 \le k \le n)
\]
are derivable in $S$, where $h_1, \cdots, h_s$ is a subsequence of
$1, \cdots, k$ and the $\mathfrak{B}_m$ are of the form $c_1, \cdots, c_m
\to \mathfrak{A}_n$ and the $\bar{\mathfrak{X}}_k$ are the formulae
$(c_1, \cdots, c_n \to \bar{\mathfrak{X}}_k) \to c_n$ and the
$\mathfrak{B}_{m+1}$ belong to $(c_1, \cdots, c_n, \mathfrak{X}_k) \to c_k$.

Hence the formulae
\[
c_1, \cdots, c_n, c_{k+1}, \cdots, c_n \to \mathfrak{B}_m^{[k]},
\bigl[(c_k \to \mathfrak{B}_n^{[k]}) \to c_k\bigr],
\cdots,
\bigl[(c_k \to \mathfrak{B}_n^{[k]}) \to c_k\bigr]
\]
are derivable in $S$, and hence by Lemma~1 so are the formulae

\[
c_1, \cdots, c_n, c_{k+1}, \cdots, c_n,
\bigl[(c_k \to \mathfrak{B}_n) \to c_k\bigr],
\cdots,
\bigl[(c_k \to \mathfrak{B}_n) \to c_k\bigr]
\to c_k \quad (1 \le k \le n).
\]

Taking $L_k$ to be the set of formulae belonging to $C(L_k, c_k)$ satisfying
the conditions of Lemma~3, Corollary (i).

Hence the formulae
$(c_k \to \mathfrak{B}_n^{[k]}), (c_k \to \mathfrak{B}_n^{[k]}),
\cdots, (c_n, \to \mathfrak{B}_n^{[k]}), (c_k \to \bar{\mathfrak{X}}_k^{[k]})
\to (c_k \to \bar{\mathfrak{Z}}_k^{[k]})$
and therefore also the formulae
$c_1, c_2, \cdots, c_n, \mathfrak{Y}_{h1}, \cdots, \mathfrak{Y}_{hn},
\bar{\mathfrak{X}}_k \to \bar{\mathfrak{Z}}_k$ $(1 \le k \le n)$ are
derivable in $S$. Now $\bar{\mathfrak{X}}_k = \tau(\xi_k)$,
$\mathfrak{Y}_{h1} = \tau(\eta_{h1})$ and $\bar{\mathfrak{Z}}_k = \tau(\eta_k)$
for some terms $\xi_k$, $\eta_h$ of the language of $S$-structures.

%% -----------------------------------------------------------------------
\bigskip\noindent
Hence by Lemma~5 the formula
\[
I(a \Rightarrow (b \Rightarrow (c_1 \Rightarrow \cdots \Rightarrow
  (c_n \Rightarrow (\xi_n \Rightarrow \cdots \Rightarrow
  (\xi_k \Rightarrow \xi_k \Rightarrow \cdots)
\]
$(1 \le k \le n)$

belongs to $\mathrm{C1}(\alpha \cup \beta)$.

From these we get by repeated applications of $\alpha$(i)--(v) and
$\beta$(vi)--(vii) and cuts, the formulae

\[
I(\eta_{h_1}), \cdots, I(\eta_{h_s}),\,
A(\nu_1, \cdots, \nu_{k-1}, \nu_k, \nu_{k+1}, \cdots, \nu_n)
\]
for arbitrary terms $\nu_1, \cdots, \nu_k, \nu_{k+1}, \nu_n$ in the language
of $S$-structures.
From these we get by cuts the formula
$A(\xi_1, \cdots, \xi_k, \xi_{k+1}, \cdots, \xi_n)$.
Hence $\mathfrak{A}$ belongs to $\mathrm{C1}(\alpha \cup \beta)$.

%% -----------------------------------------------------------------------
\bigskip
\noindent
$S(\mathscr{A}) \upharpoonright D = \langle D, l, A, \Rightarrow, \bigcap\rangle$

a prime formula of the form $A(\xi_1, \cdots, \xi_n)$ is true for
$\xi_1, \cdots, \xi_n \in D$ if and only if $A(\xi_1, \cdots, \xi_n)$ holds
in $S(\mathscr{A}) \upharpoonright D$. Let the mapping $\hat{\tau}$ be such
that for any $\mathfrak{A} \in \mathscr{A}$, $\hat{\tau}(\mathfrak{A})$ is the
set of structures $S(\mathscr{B})$ with $\mathscr{B} \in K$. Then from
Lemmas~7 and~8 we get our main result.

%% -----------------------------------------------------------------------
\begin{theorem}
Let $\tilde{\mathfrak{A}}$ be a classical formula of $L$ without function
symbols and containing besides $=$ occurrences only of a single $n$-ary
predicate letter $A$. Then $\tilde{\mathfrak{A}}$ is derivable in
$\mathrm{CPL}^*$ if and only if $\tau^*(\tilde{\mathfrak{A}})$ is derivable
in $S$.
\end{theorem}

Evidently the limitation of the number of predicate letters is a trivial
restriction, since a finite number of $n_i$-ary ($n_i < n$) predicate letters
$A_1, \cdots, A_k$ can be introduced in terms of a single $n$-ary one by
putting $\mathfrak{A} \mathbin{\&} \mathfrak{B} = A(\xi_1, \cdots, \xi_n) =
B(d_1, \xi_1, \cdots, \xi_n, \xi_{n+1}, \cdots, \xi_n) =
A_i(\xi_{n+1}, \cdots, \xi_n)$.

%% -----------------------------------------------------------------------
\bigskip\noindent
The proof of the counterpart of Lemma~7 proceeds analogously to that of
Lemma~7.

For the proof of the counterpart of Lemma~8 we use Kripke structures
$\langle K, R, \mathscr{A}\rangle$, where $K$ is a family of structures, $R$
a reflexive and transitive relation and $\mathscr{A} \in K$, noting that any
such Kripke structure can be extended to another
$\langle K', R', \mathscr{A}'\rangle$ such that $K'$ is the set of structures
$S(\mathscr{B})$ with $\mathscr{B} \in K$, $\mathscr{A}' = S(\mathscr{A})$
and $R'$ is defined such that $R'(S(\mathscr{B}_1), S(\mathscr{B}_2)) =
R(\mathscr{B}_1, \mathscr{B}_2)$ for any $\mathscr{B}_1, \mathscr{B}_2 \in K$.

%% -----------------------------------------------------------------------
\begin{thebibliography}{9}
\bibitem{Prawitz1965}
  Dag Prawitz,
  \textit{Natural deduction},
  Almqvist and Wiksell, Stockholm, 1965.

\bibitem{Prawitz1970}
  \rule{2em}{0.4pt},
  \textit{Some results for intuitionistic logic with second order
  quantification rules},
  \textit{Intuitionism and proof theory} (editors A.~Kino, J.~Myhill and
  R.~E.~Vesley),
  North-Holland, Amsterdam, 1970, pp.~259--270.

\bibitem{Schuttif1950}
  Kurt Sch\"{u}tte,
  \textit{Schlussweisenkalk\"{u}le der Pr\"{a}dikatenlogik},
  \textbf{Mathematische Annalen}, vol.~22 (1950), pp.~47--65.
\end{thebibliography}

\bigskip
\begin{center}
  \textsc{Universiteit van Amsterdam}\\
  \textsc{Amsterdam, Netherlands}
\end{center}

\end{document}
